\documentclass[10pt,a4paper]{amsart}
\usepackage[margin=19mm]{geometry}
\usepackage{amsmath,amssymb,mathtools}
\usepackage[scr=rsfs,cal=euler]{mathalfa}
\usepackage{xfrac}
\usepackage[unicode,hidelinks,bookmarks=false]{hyperref}
\newcommand{\ie}{i.e.\ }
\newcommand{\cf}{cf.\ }
\newcommand{\resp}{respectively\ }
\newcommand{\Subsection}{\subsection}
\providecommand{\nobreakdash}{\nobreak\hbox{-}}
\setlength{\emergencystretch}{2em}
\multlinegap0pt
\usepackage{bm}
\usepackage{bbm}
\usepackage{dsfont}
\usepackage{xspace}
\usepackage{cancel}
\usepackage{mathtools}
\usepackage{amsthm}
\makeatletter
\@ifundefined{theorem}{\newtheorem{theorem}{Theorem}}{}
\@ifundefined{lemma}{\newtheorem{lemma}[theorem]{Lemma}}{}
\makeatother
\newcommand{\N}{\mathbb{N}}
\title[Flat degenerations of flag supermanifolds for basic Lie supe]{Flat degenerations of flag supermanifolds for basic Lie superalgebras}
\author{Ibrahim Ahmad}
\date{Source: arXiv:2504.19946v1 (CC BY 4.0)}
\begin{document}
\maketitle

\noindent\emph{Source and licence.} Flat degenerations of flag supermanifolds for basic Lie superalgebras. Version arXiv:2504.19946v1 (CC BY 4.0), distributed under the Creative Commons Attribution 4.0 International licence, as recorded in the arXiv licence metadata for this exact version and shown on the abstract page https://arxiv.org/abs/2504.19946v1. Full rights notice: licenses/arxiv-2504.19946-CC-BY-4.0.txt.\par\smallskip

\noindent\textit{Editorial scope.} Counted unit: the Lemma whose source label is \texttt{lemma-tensor-action-sign} in the article (source0.tex lines 437--450, source label \texttt{lemma-tensor-action-sign}), delivered together with the article's own complete original proof of it (source0.tex lines 452--511, one \texttt{proof} environment of 60 lines). No mathematics is written, repaired or re-derived: statement and proof are the article's own text line for line. Required context retained uncounted: none. Every definition, display and label of those lines is retained unchanged. Omitted from this package and not redistributed: 1--436, 451--451, 512--1079. Nothing inside a retained excerpt is omitted or altered: every retained body is delivered exactly as the article's lines are written, so no omission receipt of any kind accompanies this package. Citations: the retained excerpts carry no citation command at all, so no bibliography entry of the article is transcribed and no reference is redistributed. Reference adaptation: none was needed and none was made; every reference inside the delivered unit resolves inside it, so no unresolved reference or citation is produced. Printed slips kept literal: no printed word, symbol, hypothesis or number of the counted statement or of its proof was altered. Layout and class: the article's own class is \texttt{amsart} with packages including geometry, amsthm, amsmath, amssymb, amscd, latexsym, multicol, verbatim, enumitem, graphicx, xy, color, stmaryrd, tikz; the wrapper is the lane's pinned \texttt{amsart} editorial wrapper at 10pt on A4, which restates the theorem-like environments the retained text needs and adds the article's own macro definitions that the retained text needs (\texttt{\textbackslash N}), with extra wrapper packages (bm, bbm, dsfont, xspace, cancel, mathtools). The delivered numbering is the unit's own. Assets: no retained span contains a figure environment, a graphics-inclusion command, a PDF-page inclusion, a table or a TikZ picture; no image file is copied into this package, the package declares no asset dependency and no image file is redistributed with it. Licence: version arXiv:2504.19946v1 of this article is distributed under the Creative Commons Attribution 4.0 International licence, as recorded in the arXiv licence metadata for this exact version and shown on the abstract page https://arxiv.org/abs/2504.19946v1. A full rights notice is delivered at licenses/arxiv-2504.19946-CC-BY-4.0.txt. Characters: every retained excerpt is pure ASCII, so the delivered engine, XeLaTeX with its default Latin Modern text font, reports no missing character.
\begin{lemma}\label{lemma-tensor-action-sign}
    Let $k,k'\in\N$ and let $v_{k\lambda}$ and $v_{k'\lambda}$ be the highest weight vectors of $K_{\mathfrak{b}}(k\lambda)$ and $K_{\mathfrak{b}}(k'\lambda)$.
    Then we have
    \begin{equation*}
        f^{(I,\mathbf{m})}.(v_{k\lambda}\otimes v_{k'\lambda})=\sum\limits_{(J',\mathbf{m}')+(J'',\mathbf{m}'')=(I,\mathbf{m})}(-1)^{K_{J',J''}}\prod\limits_{i=1}^{n}\begin{pmatrix}
            m_i\\
            m'_i
        \end{pmatrix}f^{(J',\mathbf{m}')}v_{k\lambda}\otimes f^{(J'',\mathbf{m}'')}v_{k'\lambda}
    \end{equation*}
    where $K_{J',J''}$ is the sum
    \begin{equation*}
        K_{J',J''}=\sum^q_{j=1}\sum^q_{i=j+1} J'_j\cdot J''_i
    \end{equation*}
\end{lemma}

\begin{proof}
    We are going to prove this via induction on the sum $|I|+|\mathbf{m}|$.
    There is nothing to show for monomials consisting of one even or one odd basis element.
    Write now
    \begin{equation*}
        f^{(I,\mathbf{m})}=f_{n+q}^{m_{n+q}}\cdots f_{1}^{m_{1}}
    \end{equation*}
    and let $1\leq t\leq n+q$ be maximal such that $m_{t}\neq0$.
    We are now going to distinguish between $f_{t}$ being odd or even.

    If $f_{t}$ is even, we first recall the formula for the action of $f_{t}^{m_{t}}$ on arbitrary elementary tensors being
    \begin{equation*}
        f_{t}^{m_{t}}(v\otimes w)=\sum_{r=0}^{m_{t}}\begin{pmatrix}
            m_{t}\\
            r
        \end{pmatrix}
        f^{r}v\otimes f^{m_{t}-r}w
    \end{equation*}
    Now, refer to $1\leq t_e\leq n$ as the index in the multi-index $\mathbf{m}=(m_1,\dots,m_n)$ indicating $m_t$ in $(I,\mathbf{m})$.
    Hence, we also get, where $e_{t}$ stands for the $t$-th standard vector,
    \begin{align*}
        f^{(I,\mathbf{m})}(v_{k\lambda}\otimes v_{k'\lambda})=&f_{t}^{m_{t}}(f^{(I,\mathbf{m}-m_{t}e_{t_{e}})}(v_{k\lambda}\otimes v_{k'\lambda}))\\
        =&f_{t}^{m_{t}}\left(\sum\limits_{(J',\mathbf{m}')+(J'',\mathbf{m}'')=(I,\mathbf{m}-m_{t}e_{t})}(-1)^{K_{J',J''}}\prod\limits_{i=1}^{t_{e}-1}\begin{pmatrix}
            m_i\\
            m'_i
        \end{pmatrix}f^{(J',\mathbf{m}')}v_{k\lambda}\otimes f^{(J'',\mathbf{m}'')}v_{k'\lambda}\right)\\
        =&\sum\limits_{\substack{(J',\mathbf{m}')+(J'',\mathbf{m}'')=(I,\mathbf{m}-m_{t}e_{t_{e}})\\r=0,\dots,m_{t_{e}}}}(-1)^{K_{J',J''}}\begin{pmatrix}
            m_t\\
            r
        \end{pmatrix}\\
        &\cdot\prod\limits_{i=1}^{t_{e}-1}\begin{pmatrix}
            m_i\\
            m'_i
        \end{pmatrix}f^{(J',\mathbf{m}'+re_{t_{e}})}v_{k\lambda}\otimes f^{(J'',\mathbf{m}''+(m_{t}-r)e_{t_{e}})}v_{k'\lambda}
    \end{align*}
    which, after reindexing, is exactly the result needed.

    If $f_{t}$ is odd, we can only have $m_t=1$. 
    Again, referring to $1\leq t_{e}\leq q$ as the index in the multi-index $I=(m_{i_{q}},\dots,m_{i_{1}})$ indicating $m_{t}$ in $(I,\mathbf{m})$ yields us
    \begin{align*}
        f^{(I,\mathbf{m})}(v_{k\lambda}\otimes v_{k'\lambda})=&f_{t}(f^{(I-e_{t_{e}},\mathbf{m})}(v_{k\lambda}\otimes v_{k'\lambda}))\\
        =&f_{t}\left(\sum\limits_{(J',\mathbf{m}')+(J'',\mathbf{m}'')=(I-e_{t_{e}},\mathbf{m})}(-1)^{K_{J',J''}}\prod\limits_{i=1}^{n}\begin{pmatrix}
            m_i\\
            m'_i
        \end{pmatrix}f^{(J',\mathbf{m}')}v_{k\lambda}\otimes f^{(J'',\mathbf{m}'')}v_{k'\lambda}\right)\\
        =&\sum\limits_{(J',\mathbf{m}')+(J'',\mathbf{m}'')=(I-e_{t_{e}},\mathbf{m})}(-1)^{K_{J',J''}}\prod\limits_{i=1}^{n}\begin{pmatrix}
            m_i\\
            m'_i
        \end{pmatrix}f^{(J'+e_{t_{e}},\mathbf{m}')}v_{k\lambda}\otimes f^{(J'',\mathbf{m}'')}v_{k'\lambda}\\
        +&\sum\limits_{(J',\mathbf{m}')+(J'',\mathbf{m}'')=(I-e_{t_{e}},\mathbf{m})}(-1)^{K_{J',J''}+|J'|}\prod\limits_{i=1}^{n}\begin{pmatrix}
            m_i\\
            m'_i
        \end{pmatrix}f^{(J',\mathbf{m}')}v_{k\lambda}\otimes f^{(J''+e_{t_{e}},\mathbf{m}'')}v_{k'\lambda}
    \end{align*}
    and now, we remark that if $J'+J''=I-e_{t_{e}}$, then
    \begin{equation*}
        K_{J',(J''+e_{t_{{e}}})}=\sum^q_{j=1}\sum^q_{i=j+1} J'_j\cdot (J''+e_{t_{e}})_i=\sum^q_{j=1}\sum^{t_{e}-1}_{i=j+1} J'_j\cdot J''_i+\sum^q_{j=1}J'_j=K_{J',J''}+|J'|
    \end{equation*}
    hence proving the claim.
\end{proof}

\end{document}
